\section{Agentic MLOps}
\label{sec:agentic-mlops}

The term \textit{MLOps} (Machine Learning Operations) describes the set of practices that aims to deploy and maintain machine learning models in production reliably and efficiently. While traditional MLOps focus on cloud-centric pipelines, the emergence of LLM-based agents has introduced a new paradigm: \textit{Agentic MLOps}.

\subsection{Agents as Orchestrators}
In an agentic MLOps framework, the traditional linear pipeline is replaced by a multi-agent system capable of reasoning about the state of the project \cite{su2025daao, llmhas2025survey}. Using frameworks like LangGraph, the system can:
\begin{itemize}
    \item \textbf{Decompose Complex Tasks}: Break down a high-level request (e.g., "Optimize this model for an STM32F4 with 128KB RAM") into sub-tasks like quantization, memory analysis, and code generation.
    \item \textbf{Iterate Based on Feedback}: If the memory analysis indicates a model is too large, the agent can autonomously trigger a "Neural Network Customization" process to adapt the architecture.
\end{itemize}

\subsection{Neural Network Customization}
A key component of automated MLOps for Edge AI is the ability to programmatically modify model architectures. "Neural Network Customization" refers to the automated replacement of layers, adjustment of input shapes, or modification of operations that are incompatible with specific hardware kernels. The orchestrator performs these transformations by utilizing the code-generation and reasoning capabilities of LLMs while maintaining the functional intent of the original model.

This automation addresses the "plumbing" problem, the friction between heterogeneous tools and languages, treating hardware constraints as first-class citizens in the AI optimization loop rather than late-stage integration concerns \cite{xue2025improve}.
